\section{Introduction}
\label{sec:intro}

Serving is entering an agentic regime. A coding agent's request
stream is unlike chat traffic in three ways that matter for the
serving substrate. First, its context is large and grows: on the
393-session corpus used throughout this paper, sessions open with a
$\sim$52k-token system prompt, add a few thousand tokens per tool-call
turn, and run to a 249k-token compaction ceiling, for a mean
re-prefilled context of $\sim$117k tokens. Second, its arrival process is
completion-coupled at machine pace: 78\% of turns arrive within
seconds of the previous turn's completion, because a tool loop
resumes the moment its result returns. Third, its demand does not
yield: an agent whose request is delayed does not leave; it waits
and retries, so slowdown converts offered load into a deferred
reservoir rather than into abandonment. Under this regime the prefix
cache stops being an optimization. With a warm cache a turn prefills
only its few-thousand-token increment; with a cold one it re-prefills
the full context, a factor-of-27 difference in service demand on this
workload. The cache is therefore load-bearing state, and its dynamics
deserve the same stability analysis given to queues.

Those dynamics contain a feedback loop. A miss re-prefills the full
context and writes the resulting KV blocks back into the same finite
pool it missed in, evicting other sessions' prefixes; on this
workload the write amplification - cache written by a miss versus by
a hit - is 22.6. The effective service rate is thus a function of the
hit rate, and the hit rate is a function of the eviction pressure the
service process itself generates. Feedback of this shape is the
defining structure of metastable failures in distributed systems
\cite{bronson2021metastable,huang2022wild}: a trigger displaces the
system into a degraded state, and a sustaining effect keeps it there
after the trigger is gone.

This paper measures that failure at multi-replica scale, derives its
structure, and models it quantitatively. On 8-replica TP4 GB200
fleets serving a 480B-parameter model under open-loop agentic trace
replay, a 45-minute load surge at base rate $\ls = 0.022$
sessions/s drives the fleet into a cold state that persists as long
as we observe it: $h$ falls to 0.016--0.030 from 0.94 warm, TTFT p50
rises to 240\,s, and no recovery occurs over 73 to 195 post-surge
minutes across three runs, sustained by the same base load the fleet
served warm before the surge (Section~\ref{sec:anatomy}) \evid{C01}. The identical
perturbation applied across base rates yields a probability band
rather than a threshold: 0.017 sessions/s always recovers, 0.022
always collapses, and 0.020 splits across three outcome classes,
including two runs that collapse near minute 240 with no new
perturbation - so benchmark horizons shorter than roughly 300
minutes overstate stability at boundary operating points \evid{C02}\evid{C05}.
The cold state exhibits hysteresis: the fleet cannot re-warm from
cold at a load it serves comfortably warm, and recovery requires
shedding demand below the band's lower edge \evid{C06}. An absorbing
degraded state, probabilistic entry, and hysteretic exit are the
metastable-failure signature, here with full-context re-prefill as
the sustaining effect - and the deferred-reservoir demand that makes
the state reachable is precisely the agentic arrival structure
(Section~\ref{sec:closedloop}).

Two system parameters reshape the failure in ways that are, to our
knowledge, unmeasured elsewhere. First, a host-memory KV tier does
not remove the failure; it transforms it into a third regime,
restore-bound congestion, in which the fleet holds a high hit rate
while the host-to-device restore channel pins near its
congested-state throughput and the queue diverges slowly (Section~\ref{sec:tier})
\evid{C10}. Tier capacity - the one parameter the tier series varies -
determines whether that regime persists, falls through to cold
collapse, or completely recovers (Section~\ref{sec:sizeseries}) \evid{C13}; restore
bandwidth was not varied, and its effect on persistence is
unmeasured. In one matched same-day pair at the 0.020 boundary, the
tier converted a collapsing operating point into a complete
self-recovery \evid{C12}. Second, collapse probability depends on the
number of replicas a single router instance spans, at fixed
per-replica load and fixed per-replica physics. The 16-replica
collapse curve sits left of the 8-replica curve: the mixed-outcome
point moves from 0.020 to 0.0185 in capacity-normalized rate,
roughly 7.5\%, though the edge brackets bound the shift only loosely
(0 to $\sim$16\%). Every measured 16-replica collapse is fleet-total,
while independent 8-replica shard routers contained every collapse
inside the failing 8-replica partition (Section~\ref{sec:span}) \evid{C15}\evid{C16}\evid{C18}.
In one measured run ($n=1$), a router configuration carrying no
affinity signal ran warm at $h$ 0.53 versus 0.94 and collapsed after
the surge with no recovery phase \evid{C22}: routing affinity carries
both the warm margin and the recovery path.

The measurements are organized by a first-order theory (Section~\ref{sec:theory}).
The classical characteristic-time analysis of LRU caches, closed
with the miss-write feedback, yields a retention fixed point whose
write amplification (22.6 here) is the feedback strength; a second
feedback through pool occupancy and queueing delay creates the cold
branch under this workload's reuse gaps; and the resulting
closed-form estimates land within tens of percent of measurement:
the cold-state service bound within $\sim$15\% of the measured cold
throughput, the minimum-tier-capacity rule bracketing the measured
persistence boundary, and the restore-channel bound matching the
measured congested throughput. The theory also explains why outcomes
inside the band are frequencies rather than responses: both branches
are locally stable there, and entry is decided by fluctuations
relative to a separatrix whose empirical coordinate - the
post-perturbation drain depth $D$ - we measure directly (Section~\ref{sec:separatrix}).

We accompany the theory with a fluid model and a discrete-event
simulation (DES) sharing one calibration in which every constant is
measured or read from deployed configuration except a single fitted
scale, the effective tier capacity (Section~\ref{sec:models}) \evid{C25}. The DES
implements the deployed router algorithm with no coordinator term
and no replica-count term. It reproduces the band, all four
congested-regime discriminators in 5 of 5 seeds, and the direction
of the span shift emergently \evid{C26}. Both model collapse curves sit
left of the measured ones, a documented reservation bias, so the
model's quantitative claim is the shift and the outcome-class
structure, never absolute collapse probabilities \evid{C27}. The model's
recovery branch is a documented failure: it produces zero complete
recoveries where the hardware produced four at the same operating
point. We report the failure, four screened candidate mechanisms,
and the surviving suspect, rather than tuning the discrepancy away
(Section~\ref{sec:failures}) \evid{C30--C32}.

Contributions:

\begin{itemize}
\item A first-order theory of prefix-cache metastability: the retention
  fixed point with miss-write amplification, the occupancy-queueing
  feedback that creates the cold branch, closed-form break-point and
  minimum-tier-capacity estimates each anchored to a measured
  quantity, and the dimensionless groups that organize the regimes
  (Section~\ref{sec:theory}).
\item Measured absorbing collapse in multi-replica prefix-cached serving
  (fleets up to 16 replicas), with a collapse-probability band,
  three outcome classes, horizon dependence, and hysteresis
  (Section~\ref{sec:collapse}) \evid{C01--C07}.
\item A drain-depth separatrix: post-perturbation KV occupancy $D$
  separates immediate collapse from other outcomes in all 8 runs at
  the boundary point and adds signal beyond rate over 21 runs,
  presented as a mediating observable with stated confounds
  (Section~\ref{sec:separatrix}) \evid{C08--C09}.
\item The host-tier third regime, restore-bound congestion, and a
  tier-size series showing that capacity sets regime persistence
  while the pinned restore channel sets congested-state throughput
  (Section~\ref{sec:tier}) \evid{C10--C14}.
\item Router-span dependence of collapse probability, with config
  independence, refutation of load-imbalance and (for the measured
  router) coordinator-saturation mechanisms, an affinity-necessity
  ablation, shard containment, and per-replica scatter signatures
  (Section~\ref{sec:span}) \evid{C15--C24}.
\item A calibrated fluid+DES model with a deployed-router layer, the
  emergently reproduced span-shift direction, executed falsifiers
  reported regardless of outcome, and documented failures
  (Section~\ref{sec:models}) \evid{C25--C32}.
\item A model-level mitigation screen: basin-aware admission deferral,
  derived from the separatrix, prevents collapse in both protocols
  and dominates AIMD and early-rejection baselines under a
  persistent surge; plus router-policy emulations showing the
  collapse band is a property of the affinity-load feedback family
  (Section~\ref{sec:mitigation}).
\end{itemize}

Every claim in the paper carries an evidence grade (Section~\ref{sec:grading};
Appendix~\ref{sec:claims} gives the full register): measured facts, outcome
frequencies with their counts, model-supported statements, and open
problems are typographically and rhetorically separated, and no
single-run observation supports a probability, threshold, or trend.
All measurements come from one stack - one model, one corpus family,
one router implementation, GB200 TP4 fleets - and numeric thresholds
are scoped to it. The mechanism itself requires only an LRU prefix
cache, full-context re-prefill on miss, and demand that defers
rather than abandons (Section~\ref{sec:theory}), and we state which conclusions
rest on the mechanism versus on the stack (Section~\ref{sec:limitations}).
